\section*{NeurIPS Paper Checklist}

\begin{enumerate}

\item {\bf Claims}
    \item[] Question: Do the main claims made in the abstract and introduction accurately reflect the paper's contributions and scope?
    \item[] Answer: \answerYes{}
    \item[] Justification: The abstract and introduction (Section~\ref{sec:intro}) state five specific contributions, all of which are supported by experimental results in Sections~\ref{sec:results}--\ref{sec:negative}. Limitations are discussed explicitly.

\item {\bf Limitations}
    \item[] Question: Does the paper discuss the limitations of the work performed by the authors?
    \item[] Answer: \answerYes{}
    \item[] Justification: A dedicated Limitations section covers domain specificity (financial only), answer-type bias (numerical only), document granularity (whole-document retrieval), embedding-model coverage (one dense encoder), and API dependency.

\item {\bf Theory assumptions and proofs}
    \item[] Question: For each theoretical result, does the paper provide the full set of assumptions and a complete (and correct) proof?
    \item[] Answer: \answerNA{}
    \item[] Justification: This is an empirical benchmarking paper with no theoretical results or proofs.

    \item {\bf Experimental result reproducibility}
    \item[] Question: Does the paper fully disclose all the information needed to reproduce the main experimental results of the paper to the extent that it affects the main claims and/or conclusions of the paper (regardless of whether the code and data are provided or not)?
    \item[] Answer: \answerYes{}
    \item[] Justification: Section~\ref{sec:setup} describes the experimental setup. Appendix~\ref{app:hyperparams} provides complete hyperparameter tables. Appendix~\ref{app:prompts} provides all prompt templates. The dataset is publicly available. Random seed, temperature, and all configuration details are specified.


\item {\bf Open access to data and code}
    \item[] Question: Does the paper provide open access to the data and code, with sufficient instructions to faithfully reproduce the main experimental results, as described in supplemental material?
    \item[] Answer: \answerYes{}
    \item[] Justification: The benchmark code and evaluation scripts are released as an open-source repository. The dataset (T\textsuperscript{2}-RAGBench) is publicly available from the original authors.

\item {\bf Experimental setting/details}
    \item[] Question: Does the paper specify all the training and test details (e.g., data splits, hyperparameters, how they were chosen, type of optimizer) necessary to understand the results?
    \item[] Answer: \answerYes{}
    \item[] Justification: Section~\ref{sec:setup} describes infrastructure, document representation, and reproducibility measures. Appendix~\ref{app:hyperparams} provides a complete hyperparameter table for all methods. We use the standard train/test split from the original benchmark.

\item {\bf Experiment statistical significance}
    \item[] Question: Does the paper report error bars suitably and correctly defined or other appropriate information about the statistical significance of the experiments?
    \item[] Answer: \answerYes{}
    \item[] Justification: All pairwise method comparisons use paired bootstrap tests ($B = 10{,}000$) with Bonferroni correction at $p < 0.05$ (Section~\ref{sec:metrics}). Significance levels are reported alongside all main results.

\item {\bf Experiments compute resources}
    \item[] Question: For each experiment, does the paper provide sufficient information on the computer resources (type of compute workers, memory, time of execution) needed to reproduce the experiments?
    \item[] Answer: \answerYes{}
    \item[] Justification: Section~\ref{sec:setup} describes the compute setup: BM25 and FAISS run locally on Apple Silicon; embedding, LLM, and reranking calls use Azure AI Foundry. No GPU training is required to reproduce the main results.

\item {\bf Code of ethics}
    \item[] Question: Does the research conducted in the paper conform, in every respect, with the NeurIPS Code of Ethics \url{https://neurips.cc/public/EthicsGuidelines}?
    \item[] Answer: \answerYes{}
    \item[] Justification: The research uses publicly available data (SEC filings), does not involve human subjects, and poses no foreseeable ethical concerns. An Ethics Statement is included.

\item {\bf Broader impacts}
    \item[] Question: Does the paper discuss both potential positive societal impacts and negative societal impacts of the work performed?
    \item[] Answer: \answerYes{}
    \item[] Justification: The Ethics Statement discusses both sides: the positive externality of cheaper, more accurate retrieval in RAG systems over regulatory and financial documents, and the negative externality that stronger retrieval on SEC filings could accelerate information-asymmetry exploits (e.g., algorithmic earnings extraction). It notes that the data is already public, so the release does not change access, only speed.

\item {\bf Safeguards}
    \item[] Question: Does the paper describe safeguards that have been put in place for responsible release of data or models that have a high risk for misuse (e.g., pre-trained language models, image generators, or scraped datasets)?
    \item[] Answer: \answerNA{}
    \item[] Justification: This paper does not release new models or datasets. The released code is an evaluation framework with no misuse risk.

\item {\bf Licenses for existing assets}
    \item[] Question: Are the creators or original owners of assets (e.g., code, data, models), used in the paper, properly credited and are the license and terms of use explicitly mentioned and properly respected?
    \item[] Answer: \answerYes{}
    \item[] Justification: All datasets (FinQA, ConvFinQA, TAT-DQA, T\textsuperscript{2}-RAGBench) and models are cited with their original papers. The benchmark is used under its original license terms.

\item {\bf New assets}
    \item[] Question: Are new assets introduced in the paper well documented and is the documentation provided alongside the assets?
    \item[] Answer: \answerYes{}
    \item[] Justification: The evaluation framework code is released with documentation, configuration files, and instructions for reproducing all experiments.

\item {\bf Crowdsourcing and research with human subjects}
    \item[] Question: For crowdsourcing experiments and research with human subjects, does the paper include the full text of instructions given to participants and screenshots, if applicable, as well as details about compensation (if any)?
    \item[] Answer: \answerNA{}
    \item[] Justification: This paper does not involve crowdsourcing or research with human subjects.

\item {\bf Institutional review board (IRB) approvals or equivalent for research with human subjects}
    \item[] Question: Does the paper describe potential risks incurred by study participants, whether such risks were disclosed to the subjects, and whether Institutional Review Board (IRB) approvals (or an equivalent approval/review based on the requirements of your country or institution) were obtained?
    \item[] Answer: \answerNA{}
    \item[] Justification: This paper does not involve human subjects research.

\item {\bf Declaration of LLM usage}
    \item[] Question: Does the paper describe the usage of LLMs if it is an important, original, or non-standard component of the core methods in this research? Note that if the LLM is used only for writing, editing, or formatting purposes and does \emph{not} impact the core methodology, scientific rigor, or originality of the research, declaration is not required.
    \item[] Answer: \answerYes{}
    \item[] Justification: LLMs are core components of several retrieval methods (HyDE, Multi-Query, CRAG, Contextual Retrieval) and the generation evaluation stage. All LLM usage is documented in Sections~\ref{sec:methods}--\ref{sec:generation} with model names, temperatures, and exact prompts in Appendix~\ref{app:prompts}.

\end{enumerate}
